% template d'article LaTeX créé par Max De Wilde (UNIGE-MATIM)
% mis à jour le 2022/04/20 par Simon Hengchen (UNIGE-MATIME)
% contact : Max.DeWilde@unige.ch; Simon.Hengchen@unige.ch

\documentclass[a4paper,12pt]{article} % ce document est un article sur une feuille A4, police taille 11

\usepackage[T1]{fontenc} % compatible avec les accents
\usepackage[french]{babel} % rédigé en français
\usepackage{palatino}
\usepackage{url} % gestion des liens (on peut aussi utiliser hyperref)
\usepackage{graphicx} % gestion des images
\usepackage{array} % gestion des tableaux
\usepackage{newunicodechar} % ajouter de nouveaux caractères, ici le point médian
\newunicodechar{·}{\textperiodcentered}
\usepackage{booktabs} % permet de faire des tableaux plus lisibles

% insérez ici d'autres extensions avec la commande \usepackage[options]{nom de l'extension}


\title{Mon titre} % le titre du document
\author{Nom de l'auteur·rice} % son auteur
\date{Date du jour} % la date

\begin{document} % début du corps du texte
\maketitle % affiche le titre, l'auteur et la date

\section{Introduction} % section 1
Le texte de l'introduction... \cite{van2016introduction} % une citation

\section{État de l'art} % section 2
\subsection{Exemple de sous-section}
Le texte de l'état de l'art... \cite[p. 123]{tahmasebi2019strengths} % une citation avec n° de page
Nous pouvons également citer plusieurs ouvrages en une fois: \cite{hill2019quantifying,verborgh2013using,de2017semantic}.

\begin{figure}[h] % insère une figure ici (h = "here")
  \centering % centre la figure
  \includegraphics[scale=1]{unige-FTI} % insère une image en taille réelle
  % l'extension n'est pas précisée pour éviter des problèmes de compilation
  % plus d'info ici : http://fr.wikibooks.org/wiki/LaTeX/Inclure_des_images
  \caption{Logo de la FTI à l'UNIGE} % nom de l'image
\end{figure}

\subsection{Autre exemple de sous-section}

\section{Analyse} % section 3
Le texte de l'analyse...\footnote{\url{https://www.unige.ch/fti/fr/}, date de dernière consultation: 2022/04/20.} % \footnote crée une note de bas de page, \url permet de formatter des URLS



\begin{table}[ht] % insère un tableau ici
\centering
\begin{tabular}{lccc} % insère un tableau avec 5 colonnes, une alignée à gauche (l) et quatre au centre (c). Aligner à droite : r
\toprule % ligne horizontale en haut du tableau 
\textbf{ Mnémonique(s)/nom } & \textbf{ Type de formation } & \textbf{ Nombre de crédits } \\ % contenus des cellules séparés par des & et \\ en fin de ligne
Communication multilingue               & Bachelor               & 180 \\
Interprétation de conférence                & Master               &  90   \\
MATIM              & Master              & 90   \\
MATS; MATT; MATCOM    & Master              & 120 \\
\bottomrule
\end{tabular}
\caption{Enseignements offerts à la FTI.}
\label{tab:enseignements} % \label{nom} pour nommer un objet, \ref{nom} pour y faire référence dans le texte
\end{table}

Comme illustré dans le tableau \ref{tab:enseignements}... % nous faisons référence au tableau {tab:enseignements}

\section{Conclusion} % section 4
Le texte de la conclusion... \cite{hengchen2015extraction}

\bibliographystyle{plain} % paramètre l'affichage de la bibliographie
\bibliography{biblio} % indique dans quel fichier .bib se trouve la biblio

\end{document} % fin du corps du texte